\uuid{ZVRW}
\exo7id{4231}
\titre{exo7 4231}
\auteur{quercia}
\organisation{exo7}
\datecreate{2010-03-12}
\isIndication{false}
\isCorrection{true}
\chapitre{Calcul d'intégrales}
\sousChapitre{Somme de Riemann}
\cpge{mpsi-mp2i.fonctions-valeurs-reelles-complexes.c, pcsi.techniques-calcul-fonctions.c, ptsi.techniques-calcul-fonctions.c, mpsi-mp2i.nombres-reels-suites.d, pcsi.nombres-reels-suites.c, ptsi.nombres-reels-suites.c, pcsi.integration.f, ptsi.integration.f, bcpst1.fonctions-reelles-usuelles.c, bcpst1.suites-reelles, bcpst1.developpements-limites.a, mp-mpi.fonctions-vectorielles.e}
\module{Analyse}
\niveau{L1}
\difficulte{}

\contenu{
\begin{enumerate}
    \item \question{Montrer que : $\forall\ x \ge 0,\ x-\frac{x^2}2 \le \ln(1+x) \le x$.}
    \item \question{Trouver $\lim_{n\to\infty}\prod_{k=1}^n\left(1+\frac1{k^2+n^2}\right)^n$.}
\reponse{
$\exp\left(\frac\pi4\right)$.
}
\end{enumerate}
}
